% ===== §8 =====
\section{The Dialogic and the Unfinished}
\label{sec:dialogic}

\noindent
The last of the traditions is the one that has already, as the map recorded, set half a foot across the line the antithesis will complete. It is the dialogic tradition of Bakhtin, and it is placed here, at the end of the thesis, because it is less an opposing view to be overcome than a near-ally whose insight the paper will take up and radicalize. Where the four preceding traditions, for all their differences, tend to render culture as something settled---an acquirable whole, a prior web, a deep structure, a congealed order of power---Bakhtin's thought is organized around the refusal of settlement, and its governing concepts name, one after another, the ways in which meaning declines to close.

The first is \emph{dialogism}. For Bakhtin the word is never solitary and never first; every utterance is a link in a chain of utterances, shaped in advance by the words it answers and turned in advance toward the response it anticipates, so that meaning lives not in the utterance taken alone but in the relation between utterances, in the space of address and reply \citep{bakhtin1981}. An utterance is oriented, before it is even spoken, toward an addressee whose answering word it already reaches for; meaning is a matter of this orientation, this being-toward-the-other, rather than a content lodged in a sign. The second is \emph{answerability}, the theme of his earliest work: the self is constituted in and as its responsiveness to the other, answerable in the double sense of being obliged to respond and of being responsible in responding, so that subjectivity itself is a dialogic and ethical event before it is a substance. The third is \emph{heteroglossia}, the internal stratification of any living language into a multiplicity of social voices---dialects, registers, professional and class and generational tongues---that coexist in tension and contend within every utterance, so that a language is never one but always a contested plurality, and its unity is an achievement of force rather than a natural fact \citep{bakhtin1981}. And the fourth, the deepest, is \emph{unfinalizability}. Nothing conclusive, Bakhtin insists, has yet taken place in the world; the ultimate word about the world and about any self has not been and cannot be spoken; every apparently final word is penultimate, shadowed by the possibility of a response that would reopen it, and the self above all retains, against every external attempt to fix and complete it, the capacity to mean again \citep{bakhtin1984poetics}. In his reading of Dostoevsky this becomes \emph{polyphony}: a work in which a plurality of consciousnesses, each with its own validity, sound together without being merged into or subordinated to a single authorial voice, a whole that is genuinely a chorus and not a monologue in disguise.

\subsection*{Why the tradition leans toward this paper}

It should be plain how far this tradition has already moved toward the position the paper will occupy. Its conception of culture, drawn out of these concepts, is of something generated in the between of address and answer, existing only in the ongoing exchange, and constitutively unfinished; and this is, in its formal shape, close to what this paper will claim of the culture generated in intimate relation. The dialogic word that is oriented toward its answer is not far from the private language that two people build in the alternation of speaking and being understood; the unfinalizable self that keeps the power to mean again is the self this series has always defended against the closures that would reduce it; and the polyphony that lets voices sound together without merging is a model, at the scale of the novel, of what a good relation lets its two members be. Of all the traditions surveyed, this is the one whose truth the synthesis will most nearly keep, and it is a second witness, alongside the verbal force of \emph{hua}, to the thesis that culture is process before it is product.

\subsection*{What it lacks, and why it is not yet the synthesis}

But it is not yet the synthesis, and the difference is exactly what the antithesis and the meta-theoretical statement must add. Bakhtin gives the \emph{form} of unfinished, dialogic generation without the apparatus that would explain its \emph{depth} and its \emph{stakes}. Three things are missing, and each is supplied by a line this paper carries. First, there is no theory of \emph{drive}: the dialogic orientation toward the other touches, but does not systematize, the question of what in the subject reaches for the other's answer, and it is the psychoanalytic account of desire and its lack that gives the reaching its root, showing that the dialogue is not only an exchange of meanings but a circulation of desire around a constitutive emptiness. Second, there is no articulated theory of the \emph{registers}: Bakhtin's dialogism lives almost entirely in the medium of the word, the symbolic, and does not thematize the imaginary body of identification or the real of the unsymbolizable bond, which the couple's culture, this paper will argue, inhabits no less than it inhabits language. Third, there is no \emph{materialism} of the kind the paper insists on: heteroglossia registers that language is a field of social contention, and gestures thereby toward power, but it does not anchor the dialogic process in the forces and relations of production that condition which voices may be raised and which answers may be given. Bakhtin, in short, has the generativity and the unfinishedness, and lacks the drive, the registers, and the material ground that would turn a poetics of the unfinished word into an ontology of cultural emergence. The thesis therefore ends where it should, at a tradition that has named the destination without yet possessing the means to reach it. The means are the business of the antithesis, and the antithesis begins, as it must, not by adding one more view to these six but by stepping back to state the framework within which all six will be reread.
